%% kapa: exam=1 round=35 year=2024 period=2 form=A subject=accounting subject_ko=회계학 seq=26 number=66
%% source: Q-Net 감정평가사 기출문제 (cst00309) artlSeq=5213229 file=raw/1008/5213229/2교시 문제지 원본.pdf
%% conversion: kordoc 4.24.51 -> tools/gen1.py (auto); review=proofread-sonnet; flags=none
\begin{question}{66}
생물자산에 관한 설명으로 옳지 않은 것은?
\begin{choices}
  \item 어떠한 경우에도 수확시점의 수확물은 공정가치에서 처분부대원가를 뺀 금액으로 측정한다.
  \item 수확 후 조림지에 나무를 다시 심는 원가는 생물자산의 원가에 포함된다.
  \item 최초의 원가 발생 이후에 생물적 변환이 거의 일어나지 않는 경우 원가가 공정가치의 근사치가 될 수 있다.
  \item 생물자산이나 수확물을 미래 일정시점에 판매하는 계약을 체결할 때, 공정가치는 시장에 참여하는 구매자와 판매자가 거래하게 될 현행시장의 상황을 반영하기 때문에 계약가격이 공정가치의 측정에 반드시 목적적합한 것은 아니다.
  \item 생물자산이나 수확물을 유의적인 특성에 따라 분류하면 해당 자산의 공정가치 측정이 용이할 수 있을 것이다.
\end{choices}
\end{question}
\answer{2}
